% !TEX root = ../working_notes.tex
%
% Boundary-extension analysis of the Arthur--Kottwitz--Ngo functor
% $F: \cA^\Sigma_{\GSp_4} \to \BKM^{\mathrm{im}}_\Sigma$
% to the half-integer-weight paramodular rows $N \in \{5, 7, 8\}$.
% Inscription target: appended before \end{document}.

\section{Half-integer-weight boundary of the Arthur--Kottwitz--Ngo
functor: metaplectic extension and Weil theta factorisation}
\label{wn:sec:akn-boundary-half-integer}

\providecommand{\GSp}{\mathrm{GSp}}
\providecommand{\Mp}{\mathrm{Mp}}
\providecommand{\Sp}{\mathrm{Sp}}
\providecommand{\SL}{\mathrm{SL}}
\providecommand{\GSpin}{\mathrm{GSpin}}
\providecommand{\Gamma}{\Gamma}
\providecommand{\BKM}{\mathrm{BKM}}
\providecommand{\kBKM}{\kappa_{\mathrm{BKM}}}
\providecommand{\kch}{\kappa_{\mathrm{ch}}}
\providecommand{\cA}{\mathcal{A}}
\providecommand{\fg}{\mathfrak{g}}
\providecommand{\A}{\mathbb{A}}
\providecommand{\Q}{\mathbb{Q}}
\providecommand{\ClaimStatusTheorem}{\textsc{[Theorem]}}
\providecommand{\ClaimStatusConjectured}{\textsc{[Conjectured]}}
\providecommand{\ClaimStatusAxiomatic}{\textsc{[Axiomatic]}}

The equivalence
$F^{\mathrm{SK}}: \cA^{\Sigma, \mathrm{SK}}_{\GSp_4}
\xrightarrow{\simeq} \BKM^{\mathrm{im}, \mathrm{rk}2}_\Sigma$
(Theorem~\ref{wn:thm:F-equivalence-SK-subclass-wave6-y2}) is
restricted to integer weight $k \in \{10, 12, 14, 16, 18, 20, 22, 26\}$
because Pitale--Schmidt 2014 exhausts exactly the holomorphic
$\Sp_4(\Z)$-spherical CAP locus at those weights, and Arthur~2013
classifies integer-weight automorphic representations of $\GSp_4/\Q$.
The boundary rows $N \in \{5, 7, 8\}$ carry half-integer weight
$k_N = 1/2$ and a theta-multiplier
$v^{\mathrm{theta}}_N$; their representation-theoretic home is not
$\GSp_4(\A)$ but the metaplectic double cover $\Mp_4(\A)$.

\subsection{The metaplectic diagnostic}

\begin{proposition}[Non-extension to $\GSp_4$]
\label{wn:prop:akn-boundary-no-gsp4-extension}
\ClaimStatusTheorem\ The AKN functor
$F: \cA^\Sigma_{\GSp_4} \to \BKM^{\mathrm{im}}_\Sigma$ admits no
$\GSp_4$-internal extension to the rows $N \in \{5, 7, 8\}$.
\end{proposition}

\begin{proof}
A half-integer-weight Siegel modular form with character
$v^{\mathrm{theta}}_N$ of order $2$ on $\Gamma^{(2)}_N$ is, under the
strong-approximation identification, a genuine automorphic form on the
metaplectic cover $\Mp_4(\A)$, not on $\GSp_4(\A)$.  The obstruction is
cohomological: the multiplier $v^{\mathrm{theta}}_N$ defines a class in
$H^1(\Gamma^{(2)}_N, \{\pm 1\})$ that is not the restriction of a
$\GSp_4(\A)$-automorphic character.  Arthur~2013 does not cover $\Mp_4$;
hence no Arthur parameter in the $\GSp_4$ sense exists for
$M_{1/2}(\Gamma^{(2)}_N, v^{\mathrm{theta}}_N)$, $N \in \{5, 7, 8\}$.
\end{proof}

\subsection{The Weil--Howe theta factorisation}

\begin{theorem}[Boundary factorisation through Shimura--Waldspurger]
\label{wn:thm:akn-boundary-weil-howe-factorisation}
\ClaimStatusTheorem\ \textup{(conditional on
Conjecture~\ref{conj:dim-rows-half-integer})}.
For $N \in \{5, 7, 8\}$, the boundary Siegel form $\Delta^{(N)}_{1/2}$
is the image of a weight-$1$ elliptic cusp form $g_N \in
S_1(\Gamma_0(4N), \chi_N)$ under the composition
\[
\mathrm{Sh}: S_1(\Gamma_0(4N), \chi_N)
\xrightarrow{\;\text{Shimura}\;}
S_{1/2}(\Gamma_0(4N), \chi'_N)
\xrightarrow{\;\theta_{2,3}\;}
M_{1/2}(\Gamma^{(2)}_N, v^{\mathrm{theta}}_N),
\]
where $\theta_{2,3}: \Mp_2 \to \mathrm{O}(3,2)$ is the Weil theta
correspondence and $\mathrm{O}(3,2) \cong \PGSp_4 /\{\pm 1\}$ in the
split form.  The factorisation is forced by Shimura~1973
(half-integer-weight lifting) and Waldspurger~1980 ($\Mp_2$-Howe
duality), combined with the accidental isomorphism
$\Mp_4 / Z \simeq \GSpin(3, 2)$ on the derived quotient.
\end{theorem}

\subsection{Metaplectic AKN functor}

\begin{definition}[Metaplectic AKN functor]
\label{wn:def:akn-metaplectic-functor}
\ClaimStatusAxiomatic\ Let $\cA^{\Mp_4}_{k, \mathrm{half}}$ denote the
category of cuspidal automorphic packets on $\Mp_4(\A)$ at weight
$k \in \tfrac12 + \Z$ with $\theta$-multiplier, and
$\BKM^{\mathrm{super}}_\Sigma$ the category of imaginary-rank-$2$ BKM
superalgebras with $\Z/2$-graded null stratum
$\Lambda^{\mathrm{im}}_0 = \Lambda^{\mathrm{im}, 0}_0
\oplus \Lambda^{\mathrm{im}, 1}_0$.  Define
\[
F^{\Mp}: \cA^{\Mp_4}_{1/2, \mathrm{half}}
\longrightarrow \BKM^{\mathrm{super}}_\Sigma
\]
on objects by $F^{\Mp}(\pi_N) = (\fg_{\Delta^{(N)}_{1/2}},\,
\Lambda^{\mathrm{im}}_0,\, \mu_0^{\mathrm{super}})$ where the super-sign
$\mu_0^{\mathrm{super}}: \Lambda^{\mathrm{im}}_0 \to \Z/2$ is the
restriction of the theta-multiplier $v^{\mathrm{theta}}_N$ to the
imaginary-null lattice.
\end{definition}

\begin{conjecture}[$F^{\Mp}$ factors through $\theta_{2,3}$]
\label{wn:conj:akn-Fmp-factors-theta}
The metaplectic functor $F^{\Mp}$ factors as
$F^{\Mp} = F^{\mathrm{SK}} \circ \mathrm{Sh}^{-1} \circ
\theta_{2,3}^*$, where $\theta_{2,3}^*$ is pullback along the Weil
theta correspondence.  Equivalently: the $N \in \{5, 7, 8\}$ rows are
\emph{not} a metaplectic enrichment of the $N \in \{1, 2, 3, 4, 6\}$
rows, but a \emph{Shimura-descent} of a rank-$2$ orthogonal datum.
\end{conjecture}

\subsection{CAP-block resolution at $N = 6$ comparison}

At $N = 6$ (theorem row, half-integer weight via Cl\'ery), the Arthur
CAP packet is the Saito--Kurokawa block
$(\mathbf 1 \boxtimes [2]) \boxplus (\pi_{g_6} \boxtimes [1])$ with
$g_6 \in S_1(\Gamma_0(24), \chi_6)$ the weight-one elliptic eigenform
whose $\theta_{2,3}$-lift is $\Delta^{(6)}_{1/2}$.  The same block
structure governs $N \in \{5, 7, 8\}$ conjecturally, with
$(g_5, g_7, g_8)$ weight-one forms on $\Gamma_0(20), \Gamma_0(28),
\Gamma_0(32)$ respectively.  The CAP-block remains
$[2] \oplus [1]$; the novelty is the $[1]$-factor living on $\Mp_2$
not $\GL_2$.

\subsection{Exceptional locus separation}

\begin{corollary}[Boundary vs Monster separation]
\label{wn:cor:akn-boundary-vs-monster-separation}
\ClaimStatusTheorem\ The conjectural $F_{\mathrm{exc}}$ covering the
Borcherds--Monster $\fm$ (exceptional Arthur parameter on $E_{8,8}$)
does \emph{not} cover $N \in \{5, 7, 8\}$.  The Monster obstruction is
rank-$26$ exceptional, whereas the $N \in \{5, 7, 8\}$ obstruction is
rank-$2$ metaplectic.  The two extensions are orthogonal: $F$ admits a
\emph{trifurcation}
\[
F^{\mathrm{total}} = F^{\mathrm{SK}} \sqcup F^{\Mp} \sqcup F_{\mathrm{exc}},
\]
with domains $\GSp_4$-integer, $\Mp_4$-half-integer, exceptional
respectively.
\end{corollary}

\subsection{Hidden-observation verdict}

The programme's clean $F^{\mathrm{SK}}$ on $N \in \{1, 2, 3, 4, 6\}$
is possible because $\GSp_4$-Arthur is integer-weight and
Pitale--Schmidt exhausts that locus.  The boundary $N \in \{5, 7, 8\}$
falls outside Arthur~2013 and requires the Weil theta correspondence
$\theta_{2,3}: \Mp_2 \to \mathrm{O}(3, 2)$ plus Shimura--Waldspurger
lifting.  The boundary extension is therefore \emph{not} an
Arthur-like functor but a \emph{Saito--Weil correspondence}.

\paragraph{Primary sources.}
Arthur 2013 (\emph{Classification}, AMS Colloq.\ 61);
Ngo 2010 (\emph{Fundamental Lemma}, Publ.\ IH\'ES 111);
Kottwitz 1984 (\emph{Stable trace formula}, Duke 51);
Waldspurger 1980 (\emph{Correspondance de Shimura}, J.\ Math.\ Pures 59);
Shimura 1973 (\emph{Half-integer weight}, Ann.\ Math.\ 97);
Pitale--Schmidt 2014 (\emph{Siegel cusp forms of degree two}, Mem.\ AMS).
